\section{多模态智能体与当前挑战}

\subsection{从文本到视觉：多模态智能体}

前面讨论的智能体主要以\textbf{文本}（HTML DOM）作为环境观测。但在实际应用中，HTML 可能包含数万个 DOM 元素，直接输入语言模型既不高效也不自然。更自然的方式是直接使用\textbf{屏幕截图}作为输入。

这催生了两类视觉-语言模型在智能体中的应用：

\subsubsection{LLaVA：合成数据训练 VLM}

LLaVA 的思路类似 Orca（蒸馏）：

\begin{enumerate}
    \item 用图像的\textbf{文本描述}（而非图像本身）提示 GPT-4 生成问答对
    \item 联合微调图像编码器（CLIP）和文本解码器（instruction-tuned LLaMA）
    \item 最终得到能理解图像并生成文本回答的多模态模型
\end{enumerate}

\begin{figure}[H]
\centering
\includegraphics[width=0.95\textwidth]{slides-images/slide-067.jpg}
\caption{LLaVA 架构：CLIP 编码器处理图像，通过投影层连接到 LLaMA 解码器\protect\footnotemark}
\end{figure}
\footnotetext{Slides 第68页。}

\subsubsection{Pix2Struct：截图-to-HTML 预训练}

Pix2Struct 引入了更自然的预训练任务：

\begin{itemize}
    \item 将网页截图中的部分区域\textbf{遮挡}
    \item 让模型预测被遮挡区域对应的 \textbf{HTML 代码}
\end{itemize}

这种预训练方式建立了视觉表示与结构化文本之间的直接对应关系，非常适合后续用于构建多模态智能体。

\begin{figure}[H]
\centering
\includegraphics[width=0.95\textwidth]{slides-images/slide-068.jpg}
\caption{Pix2Struct 的预训练任务：遮挡网页截图中的区域，让模型生成对应的 HTML\protect\footnotemark}
\end{figure}
\footnotetext{Slides 第69页。}

\subsection{当前挑战：巨大的性能鸿沟}

\begin{figure}[H]
\centering
\includegraphics[width=0.95\textwidth]{slides-images/slide-070.jpg}
\caption{MiniWoB++ 上的``提示鸿沟''：即使加上 BAGEL 的合成数据（深色），在简单任务上也远未达到完美。浅色为零样本性能\protect\footnotemark}
\end{figure}
\footnotetext{Slides 第71页。}

当前 LLM 智能体面临的核心挑战：

\begin{enumerate}
    \item \textbf{提示鸿沟（Prompting Gap）}：不进行大量 prompt 工程和定制 few-shot 示例，即使最好的 LLM 在最简单的环境上也表现很差
    \item \textbf{长周期规划困难}：从单步动作到5--10步任务，性能急剧下降
    \item \textbf{低级错误频发}：模型会犯人类不太会犯的低级错误（如在密码框里输入邮箱地址），且\textbf{无法从错误中恢复}
    \item \textbf{人机差距巨大}：在 WebArena 等更复杂的环境中，人类成功率与模型成功率之间存在数十个百分点的差距
\end{enumerate}

\begin{warningbox}{低级错误示例}
在 WebLinx 的一个任务中，GPT-4V 需要打开 Google Translate 并用给定的凭据登录。结果模型把\textbf{邮箱地址输入到了密码框}中，之后始终无法恢复——反复尝试输入邮箱。另一个例子中，模型需要执行搜索，却将同一个搜索词\textbf{重复了三遍}才提交。这些看似``愚蠢''的错误揭示了当前模型在环境理解和错误恢复方面的根本缺陷。
\end{warningbox}

\subsection{本章小结}

多模态智能体通过直接处理视觉输入来避免 HTML 解析的复杂性，但当前所有智能体方法都面临巨大的性能鸿沟。无论是简单的 MiniWoB 还是复杂的 WebArena，模型都远未达到人类水平，尤其在长周期规划和错误恢复方面。

%% ========== 总结与延伸 ==========

